\section{Distributed Training：Rare Failure 会变成日常事件}

Frontier run 同时依赖大量 accelerators、host、collective network、dataset storage、checkpoint storage、scheduler 和 cloud quota。单个组件的故障率即使很低，规模扩大后也会频繁中断。系统目标不是幻想零故障，而是限制 failure domain、快速检测、保存科学状态并恢复有效进展。

\subsection{Workers、network、storage 与 scheduler}

\term{collective communication}（集合通信）是 all-reduce、all-gather、reduce-scatter 等多设备协同 primitive；一次慢 worker 或 network partition 会拖住同步训练。Storage 既要持续供给 token，也要接收大 checkpoint；scheduler 则维护 membership、placement 与 restart。任何依赖都需要 health signal、timeout、backpressure 和 ownership。

\lecturefigure{05-training-failure-domains.png}{Frontier training 耦合 workers、network、storage 和 scheduler 四类 failure domain。}{本地字幕 15:00--17:20；概念重绘。}

\begin{knowledgebox}{读图：先区分 hard failure 与 silent corruption}
Hard failure 包括进程退出、设备掉线、storage timeout，容易触发 restart；silent corruption 可能表现为错误数据、collective bit error、stale code 或异常梯度，更危险。监控不仅要回答“job 还在跑吗”，还要回答“它是否仍在执行同一个有效实验”。
\end{knowledgebox}

\begin{importantbox}{Useful training rate}
设理论峰值 compute 为 $F_{peak}$，有效利用率为 $u$，因故障、恢复与空转损失的比例为 $d$，则长期 useful rate 可近似写为
\[
F_{useful}=F_{peak}\,u(1-d).
\]
其中提高 kernel MFU 只改善 $u$；若 checkpoint 太慢、故障检测迟或恢复不一致导致 $d$ 很大，昂贵硬件仍无法转化为有效训练进展。
\end{importantbox}

\subsection{本章小结}

规模让 rare failure 常态化。Distributed training 必须同时优化 steady-state utilization 与 recovery tax，并用实验一致性而非“job resumed”作为恢复成功标准。

